\chapter*{译者说明}
\addcontentsline{toc}{chapter}{译者说明}
\markboth{译者说明}{译者说明}

\textbf{底本。}本书译自 D.~Gaitsgory 为《朗兰兹纲领导论》撰写的第十二章。原书为 J.~Bernstein 与 S.~Gelbart 编，\emph{An Introduction to the Langlands Program}，Birkhäuser，2004 年；本章位于原书第 269--281 页（PDF 第 271--283 页），共 13 页。作者开篇即说明，本章由他在讲习班上所作的报告整理而成，包含三节和一个附录，目的是给读者一张进入几何朗兰兹的地图，而不是完整的论证。

\textbf{编排。}译文保留原书的全部编号：节号 $1$--$4$，节内逐条编号的段落与命题（1.1、1.5、2.4 等），以及原书的两条脚注。原书中编号段与命题、引理共用一套编号，本书照此办理。原书极少给公式编号，译文同样不另编号；正文中的交叉引用一律指向原书条号，读者可与英文原文逐条对照。

\textbf{页面对照。}原书页码与本书内容的对应如下；本书重排后不再保留原书页码。

\begin{center}\small
\begin{tabular}{ll}
\hline
原书页 & 内容 \\
\hline
269 & 题记；\S1.1--1.2 \\
270 & \S1.3--1.5 \\
271 & \S1.5 的证明；\S1.6 \\
272 & \S1.7--1.8 \\
273 & \S1.8 续；\S1.9 \\
274 & \S2.1--2.4 \\
275 & \S2.4 的证明；\S2.5--2.7 \\
276 & \S2.8；\S3.1 \\
277 & \S3.1 续；\S3.2--3.4 \\
278 & \S3.4 续；\S3.5；\S4.1 \\
279 & \S4.1 续；\S4.2--4.3 \\
280 & \S4.3 续；文献 \\
281 & 文献续 \\
\hline
\end{tabular}
\end{center}

\textbf{术语。}术语以原书用法为准，并与本译丛已出的几何朗兰兹译作保持一致：adele 作「阿代尔」，local system 作「局部系统」，lisse sheaf 作「lisse 层」（lisse 不译），moduli stack 作「模叠」，spherical representation 作「球表示」，cuspidal 作「尖点的」。有两处译名需要交代。

其一，Hecke eigen-sheaf 作「Hecke 本征层」，不取「Hecke 特征层」。原因是 \S4.3 的 character sheaf 必须译作「特征层」，两者若都作「特征层」便无法区分。本译丛 SY02-11 曾用「Hecke 特征层」，此处改从「本征层」，以保篇内一致。

其二，「几何朗兰兹」依原书题名直译；「朗兰兹对应」「朗兰兹对偶群」等沿用本译丛既有译法。全部术语另见随件的《术语表》。

\textbf{译注。}正文中由灰色竖线引出的段落是译者所加，共三类：易错处的预警、原文含混处的点拨、背景与直觉的补充。译注不代替原文，也不改写原文的论证；凡涉及数学内容的补充，均注明依据。来源分三级：A 级为原书、作者本人论文、经典教材与原始论文；B 级为 nLab、Stacks Project、ProofWiki 等白名单资源，且须与 A 级交叉印证；I 级为本译丛既有的译法。凡无外部依据、仅由上下文推断者，一律标明「译者按（上下文推断，无外部依据）」。

\textbf{公式。}原书排印紧凑，公式密集。译文中的每个公式都逐符号对照原书页面图像（150 dpi 渲染）核校，不依赖文本提取件——后者的公式一律不可用。原书用单行箭头表示的投影图（\S2.5、\S3.2、\S3.3）以数学排版重排，位置与原书一致。

\textbf{原文笔误。}凡不影响数学理解的拼写错误，译文照录原文，并在文献部分以译注统一说明；凡可能影响数学理解的疑点，一律保留原文表述、另出译注，不径改原文。本次翻译共记录疑点 17 条，全部已解，详见工程目录中的疑点清单。

\vspace{1.5em}
\noindent 数译\\
2026 年 11 月
\clearpage
